For the empirical optimizer, there is a simple reason to examine every default threshold. Conditional on a realized sample and a fixed payoff rule, the survival indicators are constant between successive values of $g/H$. On such an interval, expected survival-weighted payoff is constant, so proceeds increase linearly with face. At a threshold the affected observation still survives under the weak inequality in the default rule; crossing it can remove a payment. An empirical maximum is therefore attained at one of the candidate thresholds. Equal thresholds must be treated together so that all observations satisfying equality receive the same survival treatment.

Capacity-based recovery preserves the piecewise-linear calculation. Between thresholds, the recovery contribution is constant and the surviving promised payments scale with face. Additional premia that depend only on the default probability are also constant within each such interval. This is why the same candidate method can be used for the reported recovery and collateral scenarios without assuming that their proceeds curves are smooth. For the temperature coupon, observations with zero payoff contribute no repayment and never default on that claim; they do not create a finite positive-face threshold.

These numerical properties are distinct from the statistical behavior of a maximum. Optimizing over a sample can favor faces that happen to perform well in that sample, so averaging across seeds does not make the maximum exactly unbiased. The independent-seed standard errors quantify repetition uncertainty at the chosen sample size, while the consistency argument concerns the large-sample limit. The archived per-seed results and parameters allow these two issues to be examined separately. Neither form of numerical convergence establishes that the assumed disaster process or the additional pricing rules fit the behavior of an observed sovereign issuer.
